\documentclass{ximera}





\begin{document}

\section*{4.2 Integer, Coin, and Stamp Problems}

\subsection*{New Vocabulary}
\begin{itemize}
  \item even integer
  \item odd integer
  \item consecutive integers
  \item consecutive even integers
  \item consecutive odd integers
\end{itemize}

\section*{OBJECTIVE 1 Consecutive integer problems}

Recall that the integers are the numbers $\ldots, -4, -3, -2, -1, 0, 1, 2, 3, 4, \ldots$

An \textbf{even integer} is an integer that is divisible by 2. Examples of even integers are $-8$, $0$, and $22$.

An \textbf{odd integer} is an integer that is not divisible by 2. Examples of odd integers are $-17$, $1$, and $39$.

\textbf{Consecutive integers} are integers that follow one another in order. Examples of consecutive integers are shown below (Assume the variable $n$ represents an integer.)

$$11, 12, 13$$

$$-8, -7, -6$$

$$n, n+1, n+2$$

Examples of \textbf{consecutive even integers} are shown below (Assume the variable $n$ represents an even integer.)


$$24, 26, 28$$
$$-10, -8, -6$$
$$n, n+2, n+4$$

Examples of \textbf{consecutive odd integers} are shown below. (Assume the variable $n$ represents an odd integer.)


$$19, 21, 23$$
 $$-1, 1, 3$$
$$n, n+2, n+4$$



\section*{Discuss the Concepts}
\begin{enumerate}
  \item
  \begin{enumerate}
  \item
  How can we represent three consecutive integers using one variable?
\begin{solution}
$n, n+1, n+2$
\end{solution}
\item
Can we represent the three consecutive integers using the expressions $n+1$, $n+2$, and $n+3$ ?
\begin{solution}
Yes. However, when we solve an equation for n, we must add 1 to the solution to find the first integer, 2 to the solution to find the second integer, and 3 to the solution to find the third integer.
\end{solution}
 \item
Can we use $x$ instead of $n$ to represent the integers?
\begin{solution}
Yes, we can use any variable.
\end{solution}
\end{enumerate}
\begin{enumerate}
  \item
  How can we represent three consecutive even integers using one variable?
\begin{solution}
$n, n+2, n+4$
\end{solution}
\item
Can we represent the three consecutive even integers using the expressions $n+2$, $n+4$, and $n+6$ ?
\begin{solution}
Yes. However, when we solve an equation for $n$, we must add 2 to the solution to find the first integer, 4 to the solution to find the second integer, and 6 to the solution to find the third integer.
\end{solution}
\end{enumerate}
\begin{enumerate}
  \item  How can we represent three consecutive odd integers using one variable?
\begin{solution}
$n, n+2, n+4$
\end{solution}
\item Can we represent the three consecutive odd integers using the expressions $n+1$, $n+3$, and $n+5$ ?
\begin{solution}
Yes. However, when we solve an equation for $n$, we must add 1 to the solution to find the first integer, 3 to the solution to find the second integer, and 5 to the solution to find the third integer.
\end{solution}
  \item Explain why the following problem has no solution: The sum of two consecutive odd integers is eleven. Find the integers.
\begin{solution}
The solution of the equation $n+(n+2)=11$ is 4.5, and 4.5 is not an integer.
\end{solution}
  \item A student used the equation $3(n+1)= n+(n+2)+10$ to solve the following problem: "Find three consecutive integers such that three times the second equals ten more than the sum of the first and third." The solution of the equation yielded the integers 9, 10, and 11. To check the solution, the student replaced $n$ by 9 , then by 10 , and then by 11.10 and 11 did not check as solutions, so the student concluded that the solution was not correct. Explain the error in the student's conclusion.
\begin{solution}
The solution to the equation is $n=9$. Only 9 should be substituted for $n$ in the equation. The second integer 10 is represented by $n+1(9+1)$, and the third integer 11 is represented by $n+2(9+2)$.
\end{solution}
  \end{enumerate}
\end{enumerate}

\begin{problem}
Solve: The sum of three consecutive odd integers is 51. Find the integers.

STRATEGY  for solving a consecutive integer problem
Let a variable represent one of the integers. Express each of the other integers in terms of that variable. Remember that consecutive integers differ by 1. Consecutive even or consecutive odd integers differ by 2.

\begin{solution}
First odd integer: $n$\\
Second odd integer: $n+2$\\
Third odd integer: $n+4$
Determine the relationship among the integers.


The sum of the three odd integers is 51 .

$$\begin{aligned}
n+(n+2)+(n+4) & =51 \\
3 n+6 & =51 \\
3 n & =45 \\
n & =15
\end{aligned}$$

$n+2=15+2=17 \quad$ - Substitute the value of $n$ into the variable expressions $n+4=15+4=19 \quad$ for the second and third integers.\\
The three consecutive odd integers are 15, 17, and 19 .
\end{solution}
\end{problem}

\section*{In-Class Examples (Objective 4.2.1)}

\begin{problem}
The sum of three consecutive even integers is forty-eight. Find the integers.
\begin{solution}
14, 16, 18
\end{solution}
\end{problem}

\begin{problem}
Find two consecutive even integers such that three times the first is eight more than twice the second.
\begin{solution}
12, 14
\end{solution}
\end{problem}

\begin{problem}
Three times the largest of three consecutive integers is ten more than the sum of the other two integers. Find the three integers.
\begin{solution}
5, 6, and 7
\end{solution}
\end{problem}

\begin{problem}
EXAMPLE Find three consecutive even integers such that three times the second integer is six more than the sum of the first and third integers.

\begin{solution}
Strategy
First even integer: $n$\\
Second even integer: $n+2$\\
Third even integer: $n+4$\\

Three times the second integer equals six more than the sum of the first and third integers.\\
$3(n+2)=n+(n+4)+6$\\
$3 n+6=2 n+10$\\
$n+6=10$\\
$n=4$

The first even integer is 4 .\\
$n+2=4+2=6$

Substitute the value of $n$ into the
$n+4=4+4=8$ variable expressions for the second and third integers.

The three consecutive even integers are 4, 6, and 8 .
\end{solution}
\end{problem}

\begin{problem}
PROBLEM 1 Find three consecutive integers whose sum is -12.
\begin{solution}
$-5,-4,-3$
\end{solution}
\end{problem}

\section*{Optional Student Activity}

\begin{exercise}
Twice the largest of three consecutive integers is three more than the sum of the other two integers. Find the integers.
\begin{solution}
This statement is true for all sets of three consecutive integers.
\end{solution}
\end{exercise}

\begin{exercise}
If $x=2005$ represents the smallest of 2005 consecutive odd integers, what is the representation of the 2005th of these integers?
\begin{solution}
$x+4008$
\end{solution}
\end{exercise}

\section*{OBJECTIVE 2 Coin and stamp problems}

In solving problems that deal with coins or stamps of different values, it is necessary to represent the values of the coins or stamps in the same unit of money. Frequently, the unit of money is cents.\\

The value of 3 quarters in cents is $3 \cdot 25$ or 75 cents.\\

The value of 4 nickels in cents is $4 \cdot 5$ or 20 cents.\\

The value of $d$ dimes in cents is $d \cdot 10$ or $10 d$ cents.\\

\begin{problem}
Solve: A coin bank contains $\$ 1.20$ in dimes and quarters. In all, there are nine coins in the bank. Find the number of quarters in the bank.

\section*{STRATEGY for solving a coin problem}
For each denomination of coin, write a numerical or variable expression for the number of coins, the value of the coin in cents, and the total value of the coins in cents. The results can be recorded in a table:

The total number of coins is 9 .\\
Number of quarters: $x$\\
Number of dimes: $9-x$

\begin{center}
\begin{tabular}{|l|c|c|c|l|c|}
\hline
Coin & \begin{tabular}{c}
Number \\
of coins \\
\end{tabular} & . & \begin{tabular}{c}
Value of coin \\
in cents \\
\end{tabular} & $=$\begin{tabular}{c}
Total value \\
in cents \\
\end{tabular} &  \\
\hline
Quarter & $x$ & . & 25 & $=$ & $25 x$ \\
\hline
Dime & $9-x$ & . & 10 & $=$ & $10(9-x)$ \\
\hline
\end{tabular}
\end{center}

\section*{Take Note}
Use the information given in the problem to fill in the number and value columns of the table. Fill in the total value column by multiplying the two expressions you wrote in each row.

Determine the relationship between the total values of the different coins. Use the fact that the sum of the total values of the different denominations of coins is equal to the total value of all the coins.

\begin{solution}
The sum of the total values of the different denominations of coins is equal to the total value of all the coins ( 120 cents).

$$\begin{aligned}
25 x+10(9-x) & =120 \\
25 x+90-10 x & =120 \\
15 x+90 & =120 \\
15 x & =30 \\
x & =2
\end{aligned}$$

The total value of the quarters plus the total value of the dimes equals the amount of money in the bank.

There are 2 quarters in the bank.
\end{solution}
\end{problem}

\begin{problem}
EXAMPLE A collection of stamps consists of 3\textcent{} stamps and 8\textcent{} stamps. The number of 8\textcent{} stamps is five more than three times the number of 3\textcent{} stamps. The total value of all the stamps is $\$ 1.75$. Find the number of each type of stamp in the collection.

Strategy Number of 3\textcent{} stamps: $x$\\

Number of 8\textcent{} stamps: $3 x+5$

\begin{center}
\begin{tabular}{c|c|c|c}
Stamp & Number & Value & Total value \\
\hline
$3 c$ & $x$ & 3 & $3 x$ \\
\hline
$8 c$ & $3 x+5$ & 8 & $8(3 x+5)$ \\
\hline
\end{tabular}
\end{center}

The sum of the total values of the different types of stamps equals the total value of all the stamps ( 175 cents).

\begin{solution}
$$3 x+8(3 x+5)=175$$

The total value of the 3\textcent{} stamps plus the total value of the 8 stamps equals the total value of all the stamps.

$$\begin{aligned}
& 27 x=135 \\
& x=5 \\
& 3 x+5=3(5)+5 \\
&=15+5=20
\end{aligned}$$

There are five 3\textcent{} stamps and twenty 8\textcent{} stamps in the collection.
\end{solution}
\end{problem}

\begin{problem}
PROBLEM 2 A coin bank contains nickels, dimes, and quarters. There are five times as many nickels as dimes and six more quarters than dimes. The total value of all the coins is $\$ 6.30$. Find the number of each kind of coin in the bank.
\begin{solution}
40 nickels, 8 dimes, 14 quarters
\end{solution}
\end{problem}

\section*{Objective A.2.2}
\section*{Concept Check}

\begin{question}
An accountant purchased 32 stamps for $\$ 10.68$. The purchase included 24\textcent{} stamps and 39\textcent{} stamps. Write an equation to determine how many of each type of stamp were purchased.
\begin{solution}
Using decimals:\\
$0.24 x+0.39(32-x)=10.68$ or\\
$0.24(32-x)+0.39 x=10.68$;\\
Using integers:\\
$24 x+39(32-x)=1068$ or $24(32-x)+39 x=1068$\\
(The answer is twelve 24\textcent{} stamps and twenty 39\textcent{} stamps.)
\end{solution}
\end{question}

\begin{question}
The total value of the nickels and quarters in a bank is $\$ 3.30$. There are six more quarters than nickels. Write an equation to determine the number of each type of coin in the bank.
\begin{solution}
$0.05 x+0.25(x+6)=3.30$ or $5 x+25(x+6)=330$\\
(The answer is that there are 6 nickels and 12 quarters in the bank.)
\end{solution}
\end{question}

\begin{question}
A total of 24 bills is in a cash box. Some of the bills are one-dollar bills, and the rest are five-dollar bills. The total value of the cash in the box is $\$ 80$. Write an equation to determine the number of each type of bill in the cash box.
\begin{solution}
$1 x+5(24-x)=80$ or $1(24-x)+5 x=80$\\
(The answer is that the cash box contains 10 one-dollar bills and 14 five-dollar bills.)
\end{solution}
\end{question}

\section*{In-Class Examples (Objective 4.2.2)}

\begin{problem}
A coin purse contains 18 coins in dimes and quarters. The coins have a total value of $\$ 3$. Find the number of dimes and quarters in the coin purse.
\begin{solution}
8 quarters and 10 dimes
\end{solution}
\end{problem}

\begin{problem}
A postal clerk sold some 20\textcent{} stamps and some 25\textcent{} stamps. Altogether, 14 stamps were sold for a total cost of $\$ 3.05$. How many of each type of stamp were sold?
\begin{solution}
20\textcent{} stamps: 9; 25\textcent{} stamps: 5
\end{solution}
\end{problem}

\begin{problem}
A bank teller cashed a check for $\$ 185$ using five-dollar bills and twenty-dollar bills. In all, 16 bills were handed to the customer. Find the number of five-dollar bills and the number of twenty-dollar bills handed to the customer.
\begin{solution}
9 five-dollar bills and 7 twenty-dollar bills
\end{solution}
\end{problem}

\section*{Optional Student Activity}

\begin{exercise}
A coin collector went to the bank and exchanged $\$ 40$ in bills for nickels, dimes, and quarters. She received an equal number of each denomination of coin. How many of each type did she receive?
\begin{solution}
100 coins of each type
\end{solution}
\end{exercise}

\begin{exercise}
Suppose you paid $\$ 2.31$ for some candy bars at a "20\%-to-40\%-off" sale. Each candy bar sold for a whole number of cents. How many candy bars did you buy if they regularly sell for $\$ .30$ each?
\begin{solution}
11 candy bars
\end{solution}
\end{exercise}

\section*{4. Concept Review}
Determine whether the statement is always true, sometimes true, or never true.

\begin{question}
An even integer is a multiple of 2 .
\begin{solution}
Always true
\end{solution}
\end{question}

\begin{question}
When $10 d$ is used to represent the value of $d$ dimes, 10 is written as the coefficient of $d$ because it is the number of cents in one dime.
\begin{solution}
Always true
\end{solution}
\end{question}

\begin{question}
Given the consecutive odd integers -5 and -3 , the next consecutive odd integer is -1 .
\begin{solution}
Always true
\end{solution}
\end{question}

\begin{question}
If the first of three consecutive odd integers is $n$, then the second and third consecutive odd integers are represented as $n+1$ and $n+3$.
\begin{solution}
Never true
\end{solution}
\end{question}

\begin{question}
You have a total of 20 coins in nickels and dimes. If $n$ represents the number of nickels you have, then $n-20$ represents the number of dimes you have.
\begin{solution}
Never true
\end{solution}
\end{question}

\begin{question}
You have a total of seven coins in dimes and quarters. The total value of the coins is $\$ 1.15$. If $d$ represents the number of dimes you have, then $10 d+25(7-d)=1.15$.
\begin{solution}
Never true
\end{solution}
\end{question}

\section*{4.1 Exercises}
\section*{Consecutive integer problems}

\begin{exercise}
1. Integers that follow one another in order are called ? integers.
\end{exercise}

\begin{exercise}
2. Two consecutive integers differ by $\_\_\_\_$ ? Two consecutive even integers differ by $\_\_\_\_$ ? Two consecutive odd integers differ by $\_\_\_\_$ $?$
\end{exercise}

\begin{exercise}
3. Explain how to represent three consecutive integers using only one variable.
\end{exercise}

\begin{exercise}
4. Explain why both consecutive even integers and consecutive odd integers are represented algebraically as $n, n+2, n+4, \ldots$.
\end{exercise}

\begin{exercise}
5. The sum of three consecutive integers is 54 . Find the integers.
\end{exercise}

\begin{exercise}
6. The sum of three consecutive integers is 75. Find the integers.
\end{exercise}

\begin{exercise}
7. The sum of three consecutive even integers is 84 . Find the integers.
\end{exercise}

\begin{exercise}
8. The sum of three consecutive even integers is 48 . Find the integers.
\end{exercise}

\begin{exercise}
9. The sum of three consecutive odd integers is 57. Find the integers.
\end{exercise}

\begin{exercise}
10. The sum of three consecutive odd integers is 81 . Find the integers.
\end{exercise}

\begin{exercise}
11. Find two consecutive even integers such that five times the first integer is equal to four times the second integer.
\end{exercise}

\begin{exercise}
12. Find two consecutive even integers such that six times the first integer equals three times the second integer.
\end{exercise}

\begin{exercise}
13. Nine times the first of two consecutive odd integers equals seven times the second. Find the integers.
\end{exercise}

\begin{exercise}
14. Five times the first of two consecutive odd integers is three times the second. Find the integers.
\end{exercise}

\begin{exercise}
15. Find three consecutive integers whose sum is negative twenty-four.
\end{exercise}

\begin{exercise}
16. Find three consecutive even integers whose sum is negative twelve.
\end{exercise}

\begin{exercise}
17. Three times the smallest of three consecutive even integers is two more than twice the largest. Find the integers.
\end{exercise}

\begin{exercise}
18. Twice the smallest of three consecutive odd integers is five more than the largest. Find the integers.
\end{exercise}

\begin{exercise}
19. Find three consecutive odd integers such that three times the middle integer is six more than the sum of the first and third integers.
\end{exercise}

\begin{exercise}
20. Find three consecutive odd integers such that four times the middle integer is equal to two less than the sum of the first and third integers.
\end{exercise}

\begin{exercise}
21. Which of the following could not be used to represent three consecutive even integers?\\
a. $n+2, n+4, n+6$\\
b. $n, n+4, n+6$\\
c. $n+1, n+3, n+5$\\
d. $n-2, n, n+2$
\end{exercise}

\begin{exercise}
22. If $n$ is an odd integer, which expression represents the third consecutive odd integer after $n$?\\
a. $n+6$\\
b. $n+4$\\
c. $n+3$\\
d. $n+2$
\end{exercise}

\begin{exercise}
23. A bank contains 24 coins in nickels and quarters. The coins have a total value of $\$ 4.40$. Find the number of nickels and the number of quarters in the bank.
\begin{solution}
8 nickels and 16 quarters
\end{solution}
\end{exercise}

\section*{Answers to Writing Exercises}
3. Let $n$ represent the first integer. Then the second integer is $n+1$, and the third integer is $n+2$.\\
4. Consecutive even or consecutive odd integers differ by 2, regardless of whether $n$ is odd or even.

\section*{Quick Quiz (Objective 4.2.1)}

\begin{exercise}
1. The sum of three consecutive odd integers is forty-five. Find the integers.
\begin{solution}
$13,15,17$
\end{solution}
\end{exercise}

\begin{exercise}
2. Find three consecutive even integers such that four times the second is eight less than the third.
\begin{solution}
$-4,-2,0$
\end{solution}
\end{exercise}

\section*{Coin and stamp problems}

\begin{exercise}
24. A drawer contains some 18\textcent{} stamps and some 25\textcent{} stamps. Altogether, 20 stamps are in the drawer with a total value of $\$ 4.23$. How many of each type of stamp are in the drawer?
\begin{solution}
18\textcent{} stamps: 11; 25\textcent{} stamps: 9
\end{solution}
\end{exercise}

\begin{exercise}
25. Explain how to represent the total value of $x$ nickels using the equation

Number of coins $\cdot$ Value of the coin in cents $=$ Total value in cents
\end{exercise}

\section*{Quick Quiz (Objective 4.2.2)}

\begin{exercise}
1. Suppose a coin purse contains only dimes and quarters. Explain how the equation for the total value of the coins is formed.
\end{exercise}

\begin{exercise}
2. A bank contains 32 coins in nickels and dimes. Let $n$ represent the number of nickels. Complete the following table.

\begin{center}
\begin{tabular}{|l|c|c|c|c|}
\hline
Coin & \begin{tabular}{c}
Number \\
of coins \\
\end{tabular} & \begin{tabular}{c}
Value of the \\
coin in cents \\
\end{tabular} & $=$ & \begin{tabular}{c}
Total value of the \\
coins in cents \\
\end{tabular} \\
\hline
Nickel & $n$ &  & $=$ & $?$ \\
\hline
Dime & $?$ & $?$ & $=$ & $?$ \\
\hline
\end{tabular}
\end{center}
\end{exercise}

\begin{exercise}
26. The total value of the 32 coins in Exercise 25 is $\$ 2.70$. The total value of the coins in cents is $\_\_\_\_$ Use this information and the information in the table in Exercise 25 to write an equation that can be solved to find the number of nickels in the bank: $\_\_\_\_$ $+$ $\_\_\_\_$ - $\_\_\_\_$ .
\end{exercise}

\begin{exercise}
27. a. How many pennies are in $d$ dollars?\\
b. How many nickels are in $d$ dollars?\\
c. How many quarters are in $d$ dollars?
\end{exercise}

\begin{exercise}
28. a. How many cents are $n$ 39\textcent{} stamps worth?\\
b. How many dollars are $n$ 39\textcent{} stamps worth?
\end{exercise}

\begin{exercise}
29. A bank contains 27 coins in dimes and quarters. The coins have a total value of $\$ 4.95$. Find the number of dimes and quarters in the bank.
\end{exercise}

\begin{exercise}
30. A coin purse contains 18 coins in nickels and dimes. The coins have a total value of $\$ 1.15$. Find the number of nickels and dimes in the coin purse.
\end{exercise}

\begin{exercise}
31. A business executive bought 40 stamps for $\$ 14.40$. The purchase included 39\textcent{} stamps and 24\textcent{} stamps. How many of each type of stamp were bought?
\end{exercise}

\begin{exercise}
32. A postal clerk sold some 34\textcent{} stamps and some 23\textcent{} stamps. Altogether, 15 stamps were sold for a total cost of $\$ 4.44$. How many of each type of stamp were sold?
\end{exercise}

\begin{exercise}
33. A drawer contains 29\textcent{} stamps and 3\textcent{} stamps. The number of 29\textcent{} stamps is four less than three times the number of 3\textcent{} stamps. The total value of all the stamps is $\$ 1.54$. How many 29\textcent{} stamps are in the drawer?
\end{exercise}

\begin{exercise}
34. The total value of the dimes and quarters in a bank is $\$ 6.05$. There are six more quarters than dimes. Find the number of each type of coin in the bank.
\end{exercise}

\begin{exercise}
35. A child's piggy bank contains 44 coins in quarters and dimes. The coins have a total value of $\$ 8.60$. Find the number of quarters in the bank.
\end{exercise}

\begin{exercise}
36. A coin bank contains nickels and dimes. The number of dimes is 10 less than twice the number of nickels. The total value of all the coins is $\$ 2.75$. Find the number of each type of coin in the bank.
\end{exercise}

\begin{exercise}
37. A total of 26 bills are in a cash box. Some of the bills are one-dollar bills, and the rest are five-dollar bills. The total amount of cash in the box is $\$ 50$. Find the number of each type of bill in the cash box.
\end{exercise}

\begin{exercise}
38. A bank teller cashed a check for $\$ 200$ using twenty-dollar bills and ten-dollar bills. In all, 12 bills were handed to the customer. Find the number of twenty-dollar bills and the number of ten-dollar bills.
\end{exercise}

\begin{exercise}
39. A coin bank contains pennies, nickels, and dimes. There are six times as many nickels as pennies and four times as many dimes as pennies. The total amount of money in the bank is $\$ 7.81$. Find the number of pennies in the bank.
\end{exercise}

\begin{exercise}
40. A coin bank contains pennies, nickels, and quarters. There are seven times as many nickels as pennies and three times as many quarters as pennies. The total amount of money in the bank is $\$ 5.55$. Find the number of pennies in the bank.
\end{exercise}

\begin{exercise}
41. A collection of stamps consists of 22\textcent{} stamps and 40\textcent{} stamps. The number of 22\textcent{} stamps is three more than four times the number of 40\textcent{} stamps. The total value of the stamps is $\$ 8.34$. Find the number of 22\textcent{} stamps in the collection.
\end{exercise}

\begin{exercise}
42. A collection of stamps consists of 2\textcent{} stamps, 8\textcent{} stamps, and 14\textcent{} stamps. The number of 2\textcent{} stamps is five more than twice the number of 8\textcent{} stamps. The number of 14\textcent{} stamps is three times the number of 8\textcent{} stamps. The total value of the stamps is $\$ 2.26$. Find the number of each type of stamp in the collection.
\end{exercise}

\begin{exercise}
43. A collection of stamps consists of 3\textcent{} stamps, 7\textcent{} stamps, and 12\textcent{} stamps. The number of 3\textcent{} stamps is five less than the number of 7\textcent{} stamps. The number of 12\textcent{} stamps is one-half the number of 7\textcent{} stamps. The total value of all the stamps is $\$ 2.73$. Find the number of each type of stamp in the collection.
\end{exercise}

\begin{exercise}
44. A collection of stamps consists of 2\textcent{} stamps, 5\textcent{} stamps, and 7\textcent{} stamps. There are nine more 2\textcent{} stamps than 5\textcent{} stamps and twice as many 7\textcent{} stamps as 5\textcent{} stamps. The total value of the stamps is $\$ 1.44$. Find the number of each type of stamp in the collection.
\end{exercise}

\begin{exercise}
45. A collection of stamps consists of 6\textcent{} stamps, 8\textcent{} stamps, and 15\textcent{} stamps. The number of 6\textcent{} stamps is three times the number of 8\textcent{} stamps. There are six more 15\textcent{} stamps than there are 6\textcent{} stamps. The total value of all the stamps is $\$ 5.16$. Find the number of each type of stamp.
\end{exercise}

\begin{exercise}
46. A child's piggy bank contains nickels, dimes, and quarters. There are twice as many nickels as dimes and four more quarters than nickels. The total value of all the coins is $\$ 9.40$. Find the number of each type of coin.
\end{exercise}

\section*{Answers to Writing Exercises}
23. If $x$ is the number of nickels, then the total value of $x$ nickels is $x \cdot 5=$ total value in cents\\
24. The total amount of money in the coin purse can be found by finding the total value in dimes and adding it to the total value in quarters.

\subsection*{4.2 Applying Concepts}

\begin{exercise}
47. Find four consecutive even integers whose sum is $-36,-12,-10,-8,-6$
\end{exercise}

\begin{exercise}
48. Find four consecutive odd integers whose sum is $-48 .-15,-13,-11,-9$
\end{exercise}

\begin{exercise}
49. A coin bank contains only dimes and quarters. The number of quarters in the bank is two less than twice the number of dimes. There are 34 coins in the bank. How much money is in the bank?
\begin{solution}
$\$ 6.70$
\end{solution}
\end{exercise}

\begin{exercise}
50. A postal clerk sold twenty stamps to a customer. The number of 39\textcent{} stamps purchased was two more than twice the number of 24\textcent{} stamps purchased. If the customer bought only 39\textcent{} stamps and 24\textcent{} stamps, how much money did the clerk collect from the customer?
\begin{solution}
\$6.90
\end{solution}
\end{exercise}

\begin{exercise}
51. Find three consecutive odd integers such that the sum of the first and third integers is twice the second integer.
\begin{solution}
Any three consecutive odd integers
\end{solution}
\end{exercise}

\begin{exercise}
52. Find four consecutive integers such that the sum of the first and fourth integers equals the sum of the second and third integers.
\begin{solution}
Any four consecutive integers
\end{solution}
\end{exercise}

\section*{Section A.2}
\section*{Suggested Assignment}
Exercises 1-45, odds\\
More challenging problems:\\
Exercises 47-51, odds

\section*{SECTION 4.2 Answers}
\begin{itemize}
  \item[1.] consecutive
  \item[5.] The integers are 17, 18, and 19.
  \item[7.] The integers are 26, 28, and 30.
  \item[9.] The integers are 17, 19, and 21.
  \item[11.] The integers are 8 and 10.
  \item[13.] The integers are 7 and 9.
  \item[15.] The integers are $-9, -8$, and $-7$.
  \item[17.] The integers are 10, 12, and 14.
  \item[19.] No solution
  \item[21.] b
  \item[25.] Row 1: 5; $5n$; Row 2: $32-n$; 10; $10(32-n)$
  \item[27.] a. $100d$ \quad b. $20d$ \quad c. $4d$
  \item[29.] There are 12 dimes and 15 quarters in the bank.
  \item[31.] The executive bought eight 24\textcent{} stamps and thirty-two 39\textcent{} stamps.
  \item[33.] There are five 29\textcent{} stamps in the drawer.
  \item[35.] There are 28 quarters in the bank.
  \item[37.] There are 20 one-dollar bills and 6 five-dollar bills in the cash box.
  \item[39.] There are 11 pennies in the bank.
  \item[41.] There are twenty-seven 22\textcent{} stamps in the collection.
  \item[43.] There are thirteen 3\textcent{} stamps, eighteen 7\textcent{} stamps, and nine 12\textcent{} stamps in the collection.
  \item[45.] There are eighteen 6\textcent{} stamps, six 8\textcent{} stamps, and twenty-four 15\textcent{} stamps.
  \item[47.] The integers are $-12, -10, -8$, and $-6$.
  \item[49.] There is $\$6.70$ in the bank.
  \item[51.] For any three consecutive odd integers, the sum of the first and third integers is twice the second integer.
\end{itemize}


\end{document}
